\subsection{Uniform finite bounds for a fixed semisimple embedding}

\begin{proposition}[P-F17: the fixed-embedding gate closes for semisimple branching]
\label{prop:fixedplethysm}
Fix connected semisimple complex algebraic groups $H,G$ and an algebraic
homomorphism $\rho:H\to G$. There are finite constants $L_\rho\ge1$
and $C_\rho\ge0$, depending on this fixed homomorphism, such that the
following holds for every pair of dominant integral highest weights
$\lambda,\pi$ of $G,H$. Put
\[
 f_{\lambda,\pi}(n)
 =\dim\operatorname{Hom}_H
       (V_{n\pi}(H),V_{n\lambda}(G)|_\rho).
\]
If this ray has a nonzero positive-degree value, its canonical
highest-weight multiplicity algebra has an h.s.o.p. all of whose
degrees are $L_\rho$, and its Hilbert series has a positive form with
all denominator weights $L_\rho$. Moreover
\[
 s(f_{\lambda,\pi})\le p(f_{\lambda,\pi})\le C_\rho.
 \tag{P-F17a}
\]
The largest degree in a lexicographically minimal h.s.o.p. of that
canonical algebra, the modular index of GCT VI, Section 4.1.1,
is also at most $L_\rho$.
For a ray vanishing in all positive degrees, use the zero-constituent
convention and index zero; its isolated value at $n=0$ is not called
a global quasi-polynomial. The trivial pair $(\lambda,\pi)=(0,0)$
has $f(n)=1$ and both indices zero.

Consequently, for fixed $n$ and a fixed partition $\mu$ of length at
most $n$, there is a finite constant $C_{n,\mu}$ bounding the positivity
and saturation indices of all rays
\[
 k\longmapsto
 \dim\operatorname{Hom}_{GL_n}
 \bigl(S^{k\pi}\mathbb C^n,
       S^{k\lambda}(S^\mu\mathbb C^n)\bigr),
 \tag{P-F17b}
\]
where $\lambda,\pi$ range over their admissible partition lengths.
There is likewise one finite bound on the denominator weights in
some positive rational forms, and on the canonical modular indices,
for all nonzero such rays. No polynomial
bound on these constants in $n$, $\mu$, or the group dimension, and
no uniform algorithm producing those forms, is asserted.
\end{proposition}

\begin{proof}
We verify the algebraic inputs to P-F16, including the multiplication
on its diagonals. Choose Borel subgroups of $G,H$, with unipotent
radicals $U_G,U_H$ and maximal tori $T_G,T_H$.
For a dominant weight $\alpha$, write
$\alpha^*=-w_0\alpha$ for the highest weight of the dual representation.
Let
\[
 A_G=\mathbb C[G]^{U_G},\quad
 A_H=\mathbb C[H]^{U_H},\quad
 R_\rho=(A_G\otimes A_H)^H,
 \tag{P-F17c}
\]
where $H$ acts from the left through $\rho$ on the first factor and
by its left regular action on the second. The unipotent invariants
are for the right actions, so the left and right actions commute.

The exact primary construction is GCT VI, proof of Theorem 4.4.1,
printed pp.~75--76, equations (4.7)--(4.8): $A_G,A_H$ are finitely
generated, and their left-group/right-torus decomposition is
\[
 A_G=\bigoplus_{\lambda}V_\lambda(G)
       \quad\text{with right weight }\lambda^*,\qquad
 A_H=\bigoplus_{\nu}V_\nu(H)
       \quad\text{with right weight }\nu^*.
\]
That proof applies Hilbert finite generation for reductive invariants
to $R_\rho$. The same construction works for a homomorphism with
kernel: only the algebraic $H$-action and reductivity are used.
Furthermore $A_G\otimes A_H$ embeds in $\mathbb C[G\times H]$,
which is a domain because the connected complex groups are
irreducible. Thus $R_\rho$ is a domain.

Grade $R_\rho$ by the right $T_G\times T_H$ characters. Its component
of degree $(\lambda^*,\pi)$ is exactly
\[
 (R_\rho)_{(\lambda^*,\pi)}
   =(V_\lambda(G)\otimes V_{\pi^*}(H))^H
   \simeq\operatorname{Hom}_H(V_\pi(H),V_\lambda(G)|_\rho).
 \tag{P-F17d}
\]
All degrees lie in the product of the two dominant-weight cones.
For semisimple groups this cone is pointed: the sum of the pairings
with all simple coroots is nonnegative on it and strictly positive
on each nonzero dominant weight. This supplies the functional $\ell$
in P-F16. The zero-degree component is $\mathbb C$. A finite generating
set can be replaced by its finitely many homogeneous components;
omit scalar components. Choosing lattice bases for the character
lattices gives a fixed integral generator matrix $M_\rho$.

It remains to check the Cohen--Macaulay and global quasi-polynomial
conditions for the actual ray subalgebras, rather than merely for
some rings with equal Hilbert functions. For a dominant weight
$\lambda$ of $G$, let
\[
 C_G(\lambda)=\bigoplus_{k\ge0}(A_G)_{k\lambda^*},
\]
with its external grading by $k$; for $\lambda=0$ this means
$\mathbb C[t]$ with a trivial $G$-action. When $\lambda\ne0$,
this is the Cartan section ring of the highest-weight projective
orbit in $\mathbb P(V_{\lambda^*})$. Here is a direct verification
of its multiplication. Fix a highest vector $v\in V_{\lambda^*}$.
The functions $g\mapsto\varphi(gv)$, for
$\varphi\in V_{\lambda^*}^*=V_\lambda$, form precisely
$(A_G)_{\lambda^*}$. Their products of degree $k$ lie in
$(A_G)_{k\lambda^*}=V_{k\lambda}$. The product map is $G$-equivariant
and nonzero: the $k$th power of any nonzero such function is nonzero
in the domain $\mathbb C[G]$. Its image is therefore all the
irreducible module $V_{k\lambda}$. Thus the actual multiplication
is the Cartan multiplication, and this ray ring is the homogeneous
coordinate ring of that affine highest-weight cone. Equivalently,
the pullback along $g\mapsto gv$ realizes its degree-one coordinate
functions and all their products. A nonzero highest-weight character
takes all values in $\mathbb C^*$, so the orbit closure is the
affine cone over the projective orbit.

Define $C_H(\pi^*)$ in the same way; its degree-$k$ component is
$V_{k\pi}^*$. Give the two factors opposite degree actions of
$\mathbb G_m$. For $\gamma=(\lambda^*,\pi)\ne0$, the componentwise
matrix-coefficient identifications give an isomorphism of graded
algebras
\[
 \bigoplus_{k\ge0}(R_\rho)_{k\gamma}
 \simeq
 \left((C_G(\lambda)\otimes C_H(\pi^*))^{\mathbb G_m}\right)^H
 =\bigoplus_{k\ge0}
      (V_{k\lambda}(G)\otimes V_{k\pi}(H)^*)^H.
 \tag{P-F17e}
\]
In the rightmost expression the multiplication is the tensor product
of the Cartan multiplications just checked. This proves compatibility
with multiplication, not just equality of the displayed vector spaces.
The formula remains valid when one weight is zero: its external
section ring is $\mathbb C[t]$, and the nonzero joint weight
distinguishes different diagonal degrees.

This is exactly the ring displayed explicitly in GCT VI,
Section 4.1.2, equation (4.5), printed p.~69. Its multiplication is
the one constructed in the proof of Theorem 4.1.1,
equations (4.2)--(4.3), printed p.~67, applied to the highest-weight
cone as in Section 3.5.1, printed p.~59, with stretch degree one.
The section rings of the highest-weight cones are normal and have
rational singularities. The proof then uses the opposite-degree
$\mathbb G_m$ quotient and the reductive $H$ quotient to retain
rational singularities. Consequently the algebra in (P-F17e) is
normal and Cohen--Macaulay. These are the exact algebraic assertions
of Theorem 3.4.1(c) and the construction in Theorem 4.1.1(c), not an
inference from its numerical Hilbert series. For a ray whose algebra
is different from $\mathbb C$,
Theorem 3.4.1(b), proved there using Lemma 4.1.3, supplies global
quasi-polynomiality. The constant ring arising when all positive
degrees vanish is excluded as specified in the statement. No bound
on the numerator degree from the erroneous index in display (4.4)
is used.

All hypotheses of P-F16 now hold simultaneously for the
positive-dimensional ray algebras of the one fixed $R_\rho$.
If $R_\rho\ne\mathbb C$,
choose its homogeneous generators as above, put
\[
 L_\rho=\operatorname{lcm}(1,\ldots,m_\rho\Delta_\rho),\qquad
 C_\rho=(\dim R_\rho-1)L_\rho,
\]
and apply P-F16. These choices are independent of $\lambda,\pi$.

We also justify the claim about the lexicographically minimal h.s.o.p.;
existence of one all-$L_\rho$ h.s.o.p. alone is not the comparison
needed for that claim. Fix a ray algebra $A$ of dimension $r$ and
write $L=L_\rho$. The proof of P-F16 shows that every homogeneous
generator has a power in $A_L$. Hence the ideal generated by $A_L$
has radical $A_{>0}$. After quotienting by any partial homogeneous
system of parameters, its image still has radical the positive-degree
maximal ideal of the quotient. If that quotient has positive
dimension, no positive-dimensional minimal prime can contain the
whole image of $A_L$. There are finitely many minimal primes, and
over $\mathbb C$ a finite union of proper linear subspaces does not
cover the finite-dimensional vector space $A_L$. Choose an element
avoiding the minimal primes of maximal positive dimension. Cutting
by it lowers the dimension by one: on every component of maximal
dimension it is nonzero, and the principal ideal dimension theorem
gives the drop by exactly one. Repeat until dimension zero. This
extends any partial parameter sequence by parameters all of degree
$L$. Arrange parameter degrees nondecreasingly; permuting an h.s.o.p.
does not change its parameter property. Now fix a lexicographically
minimal such degree sequence. Its first
entry is at most $L$, by existence of the all-$L$ sequence. Given
any initial part of this minimal sequence, extend the corresponding
parameters by degree-$L$ elements as above. Lexicographic minimality
forces its next entry to be at most $L$. Induction bounds every
entry by $L$, proving the claimed modular-index bound.

The zero joint weight corresponds only to $(0,0)$; its multiplicity
function is one. Its section algebra is $\mathbb C[t]$ and the
parameter $t^{L_\rho}$ has degree $L_\rho$, with positive numerator
$1+T+\cdots+T^{L_\rho-1}$. Thus it obeys the same denominator
bound and index conclusion. If $R_\rho=\mathbb C$, every nonzero
joint-weight ray has zero positive-degree multiplicities; taking
$L_\rho=1$ and $C_\rho=0$ handles that case too. This proves
the semisimple statement.

Finally fix $n,\mu$ as in (P-F17b), put
$W=S^\mu\mathbb C^n$, and write $m=\dim W$.
The representation map $\rho_\mu:GL_n\to GL(W)$ restricts to
\[
 \overline\rho_\mu:SL_n\longrightarrow SL_m.
\]
Indeed, the determinant of its restriction is an algebraic character
of the connected semisimple group $SL_n$ and is therefore trivial.
For clarity, this last fact follows also by differentiating: the
derivative of a character vanishes on commutators, while the
semisimple Lie algebra equals its commutator algebra; a character
with zero derivative on a connected group has trivial image.

If $\alpha$ is a $GL_j$ highest weight, write $\bar\alpha$ for its
$SL_j$ highest weight, obtained by forgetting a determinant twist;
restriction of $V_\alpha(GL_j)$ is the irreducible
$V_{\bar\alpha}(SL_j)$, and $\overline{k\alpha}=k\bar\alpha$.
The source of the Hom space in (P-F17b) has central character
$z^{k|\pi|}$, while its target has central character
$z^{k|\mu||\lambda|}$. If
\[
 |\pi|\ne|\mu||\lambda|,
 \tag{P-F17f}
\]
all positive-degree multiplicities vanish. If the two integers agree,
every $SL_n$-equivariant map between these representations is already
$GL_n$-equivariant, since the central scalar subgroup acts by the
same character on both sides and $GL_n$ is generated by that subgroup
and $SL_n$. Explicitly any $g\in GL_n(\mathbb C)$ can be written
as $zh$ with $h\in SL_n$, after choosing an $n$th root of $\det g$.
Hence, on the balanced rays, the multiplicities in (P-F17b) equal
\[
 \dim\operatorname{Hom}_{SL_n}
     (V_{k\bar\pi},V_{k\bar\lambda}|_{\overline\rho_\mu}).
\]
The fixed semisimple result applies and gives $C_{n,\mu}$ and a
uniform positive-form denominator bound. The cases $n=1$ or $m=1$
reduce to trivial semisimple groups and yield the already handled
constant or zero multiplicity functions.

For the canonical modular index we must also identify multiplication,
rather than infer the claim from this equality of coefficients.
Restriction of matrix-coefficient functions from $GL_j$ to $SL_j$
gives, degree by degree, an isomorphism of the Cartan section ring
$\bigoplus_{k\ge0}V_{k\alpha}(GL_j)$ with
$\bigoplus_{k\ge0}V_{k\bar\alpha}(SL_j)$ as a graded
$SL_j$-algebra. Each degree map is nonzero (choose a matrix coefficient
nonzero at the identity) and maps between the isomorphic irreducible
$SL_j$ modules; it is therefore an isomorphism. Restriction of
functions respects multiplication, so the degree maps are compatible
with all Cartan products. If $\bar\alpha=0$, the external grading
is retained, giving $\mathbb C[t]$ on the $SL_j$ side instead of
collapsing different degrees into constants.
Apply these maps to both factors of the diagonal section algebra
$\bigoplus_k V_{k\lambda}(GL_m)\otimes V_{k\pi}(GL_n)^*$.
On a balanced ray its central $GL_n$ action is trivial in every
degree. Its $GL_n$ and $SL_n$ invariant subalgebras consequently
coincide, and the compatible restriction maps identify this canonical
graded invariant algebra with its semisimple counterpart. H.s.o.p.
degrees and their lexicographic minimum are preserved by a graded
algebra isomorphism. The canonical modular-index bound therefore
passes to (P-F17b) as well. This proves the corollary.
\end{proof}

\paragraph{What has been gained.}
For fixed semisimple branching, the unspecified realization gate in
P-F16 is now verified: the same multigraded algebra contains all the
actual canonical ray algebras, with the required multiplication and
singularity properties. For plethysm both $n$ and $\mu$ are fixed;
fixing only the shape $\mu$ while allowing $n$ to grow is not a
fixed embedding and is not covered by this conclusion.
The remaining quantitative question is how these generator and
parameter bounds grow as the embedding varies. Neither the existence
argument nor its least-common-multiple bound proves the uniform
polynomial growth required in the narrow PH2/PH3 hypotheses. The
torus counterexamples to the unrestricted rank-only formulation
remain consistent with this fixed-embedding result. No historical
priority is asserted.
